% ===================================================================
% AB-05c: Das Verb gustar (gefallen/mögen)
% Abgeleitet aus LE-05c (Score 9.1/10)
% ===================================================================

\documentclass[11pt, a4paper]{article}

% ===================================================================
% SW-MODUS TOGGLE
% ===================================================================
% Setze \swmodustrue für Schwarz-Weiß-Druck
\newif\ifswmodus
\swmodusfalse    % Standard: Farbe

\usepackage[utf8]{inputenc}
\usepackage[T1]{fontenc}
\usepackage[ngerman]{babel}
\usepackage{helvet}
\renewcommand{\familydefault}{\sfdefault}

\usepackage[margin=2cm]{geometry}
\usepackage{enumitem}
\usepackage{tabularx}
\usepackage{booktabs}

\usepackage{xcolor}
\usepackage{tcolorbox}
\tcbuselibrary{skins}

\usepackage{fancyhdr}
\usepackage{lastpage}
\usepackage{needspace}

% Farben - Basis
\definecolor{spanisch-color}{HTML}{c41e3a}
\definecolor{grau-color}{HTML}{888888}
\definecolor{hellgrau-color}{HTML}{f5f5f5}
\definecolor{orange-color}{HTML}{b35c00}

% SW-Farben (druckerfreundlich)
\definecolor{sw-dark}{HTML}{000000}
\definecolor{sw-gray}{HTML}{666666}
\definecolor{sw-white}{HTML}{ffffff}

% Bedingte Farbzuweisung
\ifswmodus
    \colorlet{spanisch}{sw-dark}
    \colorlet{grau}{sw-gray}
    \colorlet{hellgrau}{sw-white}
    \colorlet{orange}{sw-dark}
\else
    \colorlet{spanisch}{spanisch-color}
    \colorlet{grau}{grau-color}
    \colorlet{hellgrau}{hellgrau-color}
    \colorlet{orange}{orange-color}
\fi

\newcommand{\fachfarbe}{spanisch}

% Variablen
\newcommand{\thema}{Das Verb gustar (gefallen/mögen)}
\newcommand{\sequenz}{05}
\newcommand{\block}{c}
\newcommand{\fach}{Spanisch}
\newcommand{\klasse}{7}
\newcommand{\maxpunkte}{20}

% Kopf-/Fußzeile
\pagestyle{fancy}
\fancyhf{}
\renewcommand{\headrulewidth}{0.4pt}
\renewcommand{\footrulewidth}{0.4pt}
\lhead{\fach{} -- Klasse \klasse}
\chead{AB-\sequenz\block}
\rhead{}
\lfoot{}
\rfoot{Seite \thepage\ von \pageref{LastPage}}

% Befehle
\newcommand{\punktebox}[1]{\hfill\fbox{\small \_/#1}}
\newcommand{\luecke}[1]{\rule{#1}{0.4pt}}
\newcommand{\es}[1]{\textcolor{\fachfarbe}{\textbf{#1}}}

% Merkkasten (SW-kompatibel)
\newtcolorbox{merkkasten}[1][]{
    colback=hellgrau,
    colframe=\fachfarbe,
    colbacktitle=white,
    coltitle=\fachfarbe,
    boxrule=1pt,
    arc=0pt,
    left=8pt, right=8pt, top=8pt, bottom=8pt,
    fonttitle=\bfseries,
    attach boxed title to top left={yshift=-2mm, xshift=5mm},
    boxed title style={boxrule=0pt, colback=white},
    title=#1
}

% Hinweis (SW-kompatibel)
\newtcolorbox{hinweis}{
    colback=white,
    colframe=orange,
    colbacktitle=white,
    coltitle=orange,
    boxrule=1pt,
    arc=0pt,
    left=5pt, right=5pt, top=5pt, bottom=5pt,
    fonttitle=\bfseries,
    attach boxed title to top left={yshift=-2mm, xshift=5mm},
    boxed title style={boxrule=0pt, colback=white},
    title=Hinweis
}

% Aufgabe
\newenvironment{aufgabe}[1]{%
    \needspace{5\baselineskip}%
    \nopagebreak[4]%
    \vspace{10pt}
    \noindent\textbf{\textcolor{\fachfarbe}{Aufgabe #1}}
    \nopagebreak[4]%
    \vspace{5pt}
    \nopagebreak[4]%
    \begin{list}{}{\leftmargin=0pt}
    \item[]
}{%
    \end{list}
}

% ===================================================================
\begin{document}

% Kopfbereich
\begin{center}
    {\Large\bfseries\textcolor{\fachfarbe}{Arbeitsblatt: \thema}}
\end{center}

\vspace{5pt}

\noindent
\begin{tabularx}{\textwidth}{|X|X|X|}
\hline
\textbf{Name:} \luecke{4cm} &
\textbf{Klasse:} \klasse &
\textbf{Datum:} \luecke{2cm} \\
\hline
\end{tabularx}

\vspace{15pt}

% ===================================================================
% MERKKASTEN: Das Verb gustar
% ===================================================================

\begin{merkkasten}[Merkkasten: Das Verb \textit{gustar} (gefallen/mögen)]

\textbf{Besonderheit:} Im Spanischen sagt man \textit{``Mir gefällt...''} statt \textit{``Ich mag...''}

\vspace{0.5em}

\textbf{Die indirekten Objektpronomen:}
\begin{center}
\begin{tabular}{|l|l|l|}
\hline
\textbf{Pronomen} & \textbf{Deutsch} & \textbf{Beispiel} \\
\hline
me & mir & Me gusta... \\
\hline
te & dir & \luecke{3cm} \\
\hline
le & ihm/ihr/Ihnen & \luecke{3cm} \\
\hline
nos & uns & \luecke{3cm} \\
\hline
os & euch & \luecke{3cm} \\
\hline
les & ihnen/Ihnen & \luecke{3cm} \\
\hline
\end{tabular}
\end{center}

\vspace{0.5em}

\textbf{Wann gusta, wann gustan?}
\begin{itemize}
    \item \es{gusta} + Infinitiv (Verb): Me gusta \textbf{jugar} al fútbol.
    \item \es{gusta} + Singular-Nomen: Me gusta \textbf{el fútbol}.
    \item \es{gustan} + Plural-Nomen: Me gustan \textbf{los deportes}.
\end{itemize}

\end{merkkasten}

\vspace{10pt}

% ===================================================================
% AUFGABE 1: gusta oder gustan? (5P)
% ===================================================================

\begin{aufgabe}{1 -- gusta oder gustan?}
Setze die richtige Form ein: \textit{gusta} oder \textit{gustan}. \punktebox{5}
\end{aufgabe}

\begin{enumerate}[label=\alph*)]
    \item Me \luecke{2cm} el fútbol.
    \item Me \luecke{2cm} los deportes.
    \item Te \luecke{2cm} leer libros.
    \item Nos \luecke{2cm} las películas.
    \item Le \luecke{2cm} la música.
\end{enumerate}

\vspace{15pt}

% ===================================================================
% AUFGABE 2: Sätze vervollständigen (4P)
% ===================================================================

\begin{aufgabe}{2 -- Sätze vervollständigen}
Ergänze die Lücken mit dem richtigen Pronomen und \textit{gusta/gustan}. \punktebox{4}
\end{aufgabe}

\begin{hinweis}
\textbf{Pronomen:} me, te, le, nos, os, les
\end{hinweis}

\vspace{0.5em}

\begin{enumerate}[label=\alph*)]
    \item ¿\luecke{1.5cm} \luecke{2cm} jugar al fútbol? (dir)
    \item A María \luecke{1.5cm} \luecke{2cm} escuchar música. (ihr)
    \item \luecke{1.5cm} \luecke{2cm} los videojuegos. (uns)
    \item ¿\luecke{1.5cm} \luecke{2cm} bailar? (euch)
\end{enumerate}

\vspace{15pt}

% ===================================================================
% AUFGABE 3: Dialog vervollständigen (5P)
% ===================================================================

\begin{aufgabe}{3 -- Dialog führen}
Vervollständige den Dialog. Nutze die Redemittel. \punktebox{5}
\end{aufgabe}

\begin{hinweis}
\textbf{Redemittel:} ¿Qué te gusta hacer? | Me gusta... | ¿Y a ti? | ¿Te gusta...? / ¿Te gustan...? | Sí, me gusta(n)... / No, no me gusta(n)...
\end{hinweis}

\vspace{0.5em}

\noindent
\textbf{A:} \es{¡Hola!} \luecke{5cm}

\textbf{B:} \luecke{6cm} \es{¿Y a ti?}

\textbf{A:} \es{A mí me gusta escuchar música.} \luecke{4cm}

\textbf{B:} \luecke{6cm}

\textbf{A:} \es{¡Genial!}

\vspace{15pt}

% ===================================================================
% AUFGABE 4: Freie Produktion (6P)
% ===================================================================

\begin{aufgabe}{4 -- Über dich schreiben}
Schreibe 5 Sätze über deine Vorlieben. Benutze verschiedene Pronomen und Aktivitäten. \punktebox{6}
\end{aufgabe}

\begin{hinweis}
\textbf{Ideen:} el fútbol, la música, los videojuegos, las películas, leer, bailar, nadar, jugar al baloncesto, ver la tele, escuchar música
\end{hinweis}

\vspace{0.5em}

\noindent
\textbf{Beispiel:} \textit{Me gusta jugar al fútbol. A mi amigo le gustan los videojuegos.}

\vspace{0.8em}
\noindent 1. \rule{\textwidth}{0.4pt}

\vspace{0.8em}
\noindent 2. \rule{\textwidth}{0.4pt}

\vspace{0.8em}
\noindent 3. \rule{\textwidth}{0.4pt}

\vspace{0.8em}
\noindent 4. \rule{\textwidth}{0.4pt}

\vspace{0.8em}
\noindent 5. \rule{\textwidth}{0.4pt}

\vspace{20pt}
\hrule
\vspace{10pt}

\begin{flushright}
\textbf{Punkte:} \luecke{1cm} / \maxpunkte
\end{flushright}

\end{document}
